\section{讨论：对实验与仿真的意义}

\subsection{主信号应如何排序}

在同一 $P_0$ 框架下，静态 UEP 签名是背景量，轴频线谱是它的周期调制量。若 $m_1=0.03$，轴频基频幅值固定比静态签名低约 $30.5\ \mathrm{dB}$。因此实验上不能只问“轴频有没有”，而应同时给出：
\begin{itemize}
    \item $P_0$ 的假设范围；
    \item $m_n$ 的调制深度；
    \item 目标距离和观测方向；
    \item 系统在轴频附近的等效噪声谱密度；
    \item 可用积分时间。
\end{itemize}

静态签名适合做空间梯度、慢变包络和低漂移测试；轴频适合做窄带谱线检测。两者不是互相替代关系，而是同一源项的直流/慢变分量和调制分量。

\subsection{距离分层}

50--100 m：这是近场/中间场估算区。该距离适合先做实验验证，因为信号强，轴频基频在 $10^{-7}\ \mathrm{V/m}$ 量级。但几何解释要谨慎，最好用多点源模型或有限元模型进一步检查。

300--500 m：这是当前参数下更适合采用远场单偶极模型的距离段。300 m 时轴频基频约 $6.63\times10^{-9}\ \mathrm{V/m}$，已经进入对低频噪声和积分能力敏感的范围。

1000 m 及以上：静态签名仍可按 $R^{-3}$ 估算，轴频基频约 $10^{-10}\ \mathrm{V/m}$ 量级。此时实验重点应放在噪声谱密度、长时间稳定性、通道漂移和窄带积分上，而不是简单断言轴频不可见。

\subsection{尾流与工频的实验处理}

尾流/流动诱导电场的局部上限 $vB_0$ 可以很大，但它不是远场偶极子模型的一部分。若实验或仿真要研究尾流，应单独建立尾流速度场、扩散长度、观测位置和地磁场方向，而不是把 $vB_0$ 直接和 1000 m 处的轴频线谱做 dB 排名。

工频泄漏对实验非常重要，但它是干扰项。实验中应至少记录空载、泵开/关、DAQ 接地变化、BNC 走线变化和传感器短接时的 $A_{50,i}$ 与谐波幅值。只有这样，50 Hz 才能从目标低频信号中分离出去。

\subsection{建议的实验顺序}

\begin{enumerate}
    \item 在 100 m 等效量级下做通道噪声谱密度和 50 Hz 干扰测量，得到 $S_E^{1/2}(f)$、$A_{50,i}$ 和 $A_{k,i}$。
    \item 用可控注入或仿真信号验证 3 Hz、6 Hz、9 Hz 轴频线谱的窄带检测能力。
    \item 用 $P_0=\{30,100,300,1000,3000\}\ \mathrm{A\cdot m}$ 扫描仿真，检查 300 m 和 1000 m 的可探测性矩阵。
    \item 若研究 50--100 m 几何效应，使用多点源模型替代单偶极远场解释。
    \item 若研究尾流，单独建立流体-电磁耦合模型，不把尾流局部上限作为远场结论。
\end{enumerate}
